\begin{homeworkProblem}[4]
  Seja $\{f_n\}_{n\in\mathbb{Z}^{+}}$ uma sequência de funções $\cali{A}$-mensuráveis a valores complexos em $X$.
  \begin{enumerate}[a)]
    \item Assumindo que $\sum_{n=1}^{\infty}\int|f_n|\,d\mu<+\infty$, mostre que a série $\sum f_n$ converge $\mu$-q.t.p. para uma função mensurável $f$ tal que
      \begin{align*}
        \int f\,d\mu=\sum_{n=1}^{\infty}\int f_n\,d\mu.
      \end{align*}
    \item Assumindo que $f_n\to f\,\mu$-q.t.p., onde $f$ é mesurável, mostre que $\int |f_n-f|\,d\mu\to 0$ se, e somente se, $\int |f_n|\,d\mu\to \int |f|\,d\mu$. 
  \end{enumerate}
  \begin{solution}
    \begin{enumerate}[a)]
    \item Note que como cada $f_n$ é mensurável, temos que dado $k\in\mathbb{Z}^{+}$ definimos $g_k:X\to\mathbb{C}$ dada por
      \begin{align*}
        g_k(x)=\sum_{i=1}^{k}f_i(x)\in M(X,\cali{A}).
      \end{align*}
      Logo, temos que $g_k\to f:=\sum_{n=1}^{\infty}f_n$, pelo que podemos afirmar que $f$ é mensuravel (pois é limite de mensuráveis).\\
      Agora, note que para cada $k\in\mathbb{Z}^{+}$ temos que pela linearidade da integral podemos obter
      \begin{align*}
        \left|\int g_k\,d\mu\right| = \left|\int\sum_{i=1}^{k}f_i\,d\mu\right|
        =\left|\sum_{i=1}^{k}\int f_i\,d\mu\right| 
        &\leq\sum_{n=1}^{\infty}\left| \int f_n\,d\mu \right|\\
        &\leq\sum_{n=1}^{\infty}\int|f_n|\,d\mu\\
        &<+\infty.
      \end{align*}
      Pelo que temos que $g_k$ é $\mu$-integrável para todo $k\in\mathbb{Z}^{+}$.
      Logo temos que
      \begin{align*}
        |g_k(x)|=\left| \sum_{i=1}^{k}f_i(x) \right|\leq\sum_{n=1}^{\infty}\left|f_n(x)\right|=g(x),\quad\text{ para todo }k\in\mathbb{Z}^{+},
      \end{align*}
      com $g$ mensurável, pois $|f_n|$ é mensurável para todo $n$ e $g$ é limite de a série  parcial de $|f_n|$ e como vimos anteriormente temos que como cada  $|f_n|$ é mensurável, então 
      \begin{align*}
        \int g\,d\mu=\int \sum_{n=1}^{\infty}|f_n|\,d\mu=\sum_{n=1}^{\infty}\int |f_n|\,d\mu.
      \end{align*}
      Portanto temos que $g$ é $\mu$-integrável, pelo que usando o teorema da convergência dominada temos que a série $\sum f_n$ converge $\mu$-q.t.p para uma função mensurável $f$ tal que
      \begin{align*}
        \int f\,d\mu=\int \lim_{n \to \infty}g_n\,d\mu=\lim_{n \to \infty}\int g_n\,d\mu=\lim_{n \to \infty}\sum_{i=1}^{n}\int f_i\,d\mu=\sum_{n=1}^{\infty}\int f_n\,d\mu.
      \end{align*}
      \item Note que se $f_n\to f,\, \mu$-q.t.p., então temos que $\lim_{n \to \infty}f_n(x)-f(x)=0$ para todo $x\in X\setminus N$ com $\mu(N)=0$.\\
        Suficiente)\\
        Note que como $|f_n|$ é mensurável,  temos que pelo lema de Fatou como $|f_n|\to |f|$ $\mu$-q.t.p., obtemos que
        \begin{align*}
          \int |f|\,d=\int \liminf |f_n|\,d\mu\leq \liminf \int |f_n|\,d\mu.
        \end{align*}
        Agora, note que como $|f_n|,|f_n-f|$ e $|f|$ são mensuráveis e $|f_n|\leq |f_n-f|+|f|$ temos que pela linearidade da integral em funções de $M^{+}(X,\cali{A})$ afirmamos que
        \begin{align*}
          \int |f_n|\,d\mu \leq \int |f_n-f|\,d\mu + \int |f|\,d\mu,\quad\text{ para todo }m\in\mathbb{Z}^{+}.
        \end{align*}
        Logo temos que como $\int |f_n-f|\,d\mu\to 0$, então
        \begin{align*}
          \limsup \int |f_n|\,d\mu \leq \limsup \int |f_n-f|\,d\mu + \int |f|\,d\mu = \int|f|\,d\mu.
        \end{align*}
        Assim temos que
        \begin{align*}
          \limsup \int |f_n|\,d\mu \leq \int |f|\,d\mu \leq \liminf \int |f_n|\,d\mu.
        \end{align*}
        Isto é que
        \begin{align*}
          \lim_{n \to \infty}\int |f_n|\,d\mu = \int |f|\,d\mu.
        \end{align*}
        Necessario)\\
        Note que se $\int |f_n|\,d\mu\to \int |f|\,d\mu$, temos que pela desigualdade triangular e linearidade da integral se cumpre que
        \begin{align*}
          \lim_{n \to \infty}\int |f_n-f|\,d\mu\leq \lim_{n \to \infty} \int |f_n|\,d\mu + \int |f|\,d\mu 
        \end{align*}
  \end{enumerate} 
  \end{solution}
\end{homeworkProblem}
